\section{Software Architecture}
\label{sec:architecture}

% --- Figure: end-to-end workflow ------------------------------------------
\begin{figure*}[t]
  \centering
  \includegraphics[width=\textwidth]{pfdha_framework.pdf}
  \caption{Workflow of the prototypal PFDHA framework with coupled source-model and FDHA-model logic-tree aggregation. The framework accepts three inputs: (1)~an INI configuration file specifying the site geometry (one or more discrete sites for hazard curves, or a bounding region for hazard maps) and calculation parameters; (2)~a NRML source-model logic-tree file that references one or more underlying NRML source models defining fault geometry, magnitude--frequency distribution, kinematics, and style of faulting, and that declares the source-model epistemic uncertainties; and (3)~a NRML FDHA-model logic-tree file specifying principal and distributed model alternatives with associated branch weights. For each end-branch, an earthquake rupture forecast is built and, for every rupture, sites are partitioned by their distance~$r$ to the principal fault: sites with $r \le r_{\mathrm{threshold}}$ contribute to the principal displacement component, and sites with $r > r_{\mathrm{threshold}}$ contribute to the distributed displacement component. Per-end-branch annual rates are aggregated by combined weight to produce the final hazard curves (one curve per site) or hazard maps as the branch-weighted arithmetic mean of annual rates, together with selected fractiles.}
  \label{fig:workflow}
\end{figure*}
% -------------------------------------------------------------------------

\subsection{Design Principles}
\label{sec:architecture:principles}

The framework is designed around three principles: modularity, reproducibility, and compatibility with OpenQuake Engine conventions. Modularity is achieved by cleanly separating four PFDHA model categories---principal surface-rupture probability, principal fault displacement, distributed surface-rupture probability, and distributed fault displacement---so that any model within a category can be substituted without modifying the calculation engine, and by handling epistemic uncertainty in the seismic source through a separate, OpenQuake-conformant source-model logic tree that is decoupled from the model-selection logic tree. Reproducibility is supported by configuration-driven calculations: the INI file, together with the referenced source-model logic-tree NRML file and the referenced FDHA-model logic-tree NRML file, specifies the inputs required to reproduce a calculation, including the source model(s) and any source-model uncertainties, the site geometry (a single point, a list of discrete sites, or a grid), the target displacement levels, the investigation time, and the FDHA model selections with their parameters. Compatibility with OpenQuake conventions ensures that the framework reads standard NRML source-model XML files and standard NRML source-model logic-tree files, parses the latter through \texttt{openquake.hazardlib.logictree.SourceModelLogicTree} so that the full set of OpenQuake source-model uncertainty types is honoured transparently, uses OpenQuake's geodetic and rupture-mesh utilities, and follows the Engine's \texttt{ContextMaker} pattern for site--rupture parameter extraction \citep{Pagani2014}. Source-model XML files prepared for fault-based PSHA can therefore be reused for PFDHA when they contain the rupture-surface information required for surface-displacement calculations, ensuring consistent assumptions about fault geometry, magnitude--frequency distributions, and slip rates across hazard analyses.


Figure~\ref{fig:workflow} provides an end-to-end view of the framework. The remainder of this section describes the major components in turn: the site--rupture context system (Section~\ref{sec:architecture:context}), the model registry and adapter pattern (Section~\ref{sec:architecture:adapter}), the hazard calculation engine (Section~\ref{sec:architecture:hazard_engine}), the model logic tree (Section~\ref{sec:architecture:logic_tree}), and the configuration interface (Section~\ref{sec:architecture:config}).

The framework is implemented in Python and organised under the \texttt{openquake.fdha} namespace, mirroring the OpenQuake Engine package layout. The code is released under the GNU Affero General Public License v3.0, the same licence used by the OpenQuake Engine, to remove licensing barriers to future integration.


\subsection{Site--Rupture Context System}
\label{sec:architecture:context}

The central data structure of the framework is the \texttt{FDHAContext} dataclass, which holds all parameters required for model evaluation as flat NumPy arrays of shape~$(N,)$, where~$N$ is the number of site--rupture pairs in the calculation. For multi-site calculations~$N$ equals the number of sites times the number of surface-rupturing events retained for each site, so a single \texttt{FDHAContext} carries the full per-site, per-rupture parameter set required for vectorised model evaluation. The context fields fall into three groups: rupture parameters (magnitude~$m$, rake, dip, depth-to-top-of-rupture~$z_{\mathrm{tor}}$, and annual occurrence rate); site parameters (site identifier and~$V_{S30}$, the latter retained for forward compatibility with non-ergodic displacement models); and distance parameters (closest distance to surface trace~$r$ in km, signed perpendicular distance~$r_x$, normalised along-strike position $x/L \in [0,1]$, and rupture trace length~$L$ in km). Two derived properties are computed on demand: \texttt{style}, which classifies each rupture as normal, reverse, or strike-slip from the rake angle following \citet{AkiRichards1980}, and \texttt{hw\_fw}, which assigns each site to the hanging wall or footwall from the sign of~$r_x$.


The \texttt{FDHAContextMaker} factory class creates contexts from OpenQuake rupture objects and a site collection. This design follows the OpenQuake Engine \texttt{ContextMaker} pattern, which is used in PSHA for ground-motion model evaluation. For each rupture passed to the factory, three operations are performed: (i)~the rupture is screened for surface expression, with ruptures whose minimum mesh depth exceeds a configurable tolerance excluded from the calculation; (ii)~all distance and location metrics are computed; and (iii)~sites whose closest distance to the rupture trace exceeds the configured \texttt{maximum\_distance} are filtered out. Distance calculations are performed in a local equirectangular projection centred at the mean trace latitude, providing kilometre-scale distances that are robust to crossings of the International Date Line. The $x/L$ ratio is computed by orthogonal projection of each site onto the polyline defined by the rupture's surface trace, with both~$x$ and~$L$ evaluated within the same projection so that endpoint values of $x/L$ are exactly~$0$ or~$1$. The factory caches distance calculators by surface-geometry hash, avoiding redundant computation when many ruptures share the same surface---as occurs when ruptures of varying magnitude float along a single fault.


\subsection{Model Registry and Adapter Pattern}
\label{sec:architecture:adapter}

Models are organised into four base classes, one per PFDHA component: principal surface-rupture probability ($P[\,sr \neq 0 \mid m\,]$), principal fault displacement ($P[\,D > D_0 \mid \cdot\,]$), distributed surface-rupture probability ($P[\,d \neq 0 \mid \cdot\,]$), and distributed fault displacement ($P[\,d > d_0 \mid \cdot\,]$). The corresponding code modules are \texttt{primary\_surf\_rup}, \texttt{primary\_surf\_displ}, \texttt{secondary\_surf\_rup}, and \texttt{secondary\_surf\_displ}; the module names follow the IAEA SSG-9 ``primary/secondary'' terminology, which is equivalent to the ``principal/distributed'' terminology used throughout this paper. Each base class defines a \texttt{get\_prob()} method as the common interface, returning exceedance probabilities given the relevant subset of context parameters. This follows the OpenQuake Engine's ground-motion intensity model (GSIM) registry pattern, where each model is a self-contained, interchangeable implementation of a common interface \citep{Pagani2014}.

Different published models accept different parameter subsets and return outputs in different shapes. For example, some models use the normalised position $x/L$ while others use only distance~$r$; some return deterministic exceedance probabilities while others return Monte~Carlo samples of shape $[N, n_{\mathrm{MC}}]$ that must be reduced to deterministic values; and some models support multiple faulting styles whereas others are calibrated for a single style. To accommodate this heterogeneity without code duplication, a \texttt{LegacyModelAdapter} layer mediates between the unified context and individual model signatures. At initialisation, the adapter inspects each model's \texttt{get\_prob()} signature using Python's \texttt{inspect} module, caches the accepted parameter set, and at evaluation time automatically filters the context so that only parameters the model expects are passed. For models returning Monte~Carlo samples, the adapter applies a configurable reduction strategy---median, mean, or a specified percentile---to produce deterministic exceedance probabilities. This adapter pattern allows models originally developed with different interfaces to coexist within a single calculation without any modification to the underlying model code.


\subsection{Hazard Calculation Engine}
\label{sec:architecture:hazard_engine}

The hazard calculation is implemented in a single unified function, \texttt{calculate\_fdha\_hazard()}, that handles point-site hazard curves, multi-site hazard curves, and gridded hazard maps. The choice between curve and map is controlled by the input \texttt{[geometry]} section: a \texttt{region} declaration produces a hazard map at the configured displacement level; otherwise a hazard-curve calculation is performed for each site declared by \texttt{sites} (single \texttt{lon~lat} pair or a comma-separated list of \texttt{lon~lat} pairs with an optional per-site depth) or by \texttt{sites\_csv} (a CSV file with required columns \texttt{lon} and \texttt{lat} and an optional \texttt{depth} column). Longitude and latitude are validated against the ranges $(-180, 360]$ and $[-90, 90]$ respectively at parse time, so out-of-range coordinates are rejected before any rupture is generated. No separate calculator class is required for the curve and map output types, and no separate calculator path is required for single-site and multi-site curves; in both cases, the \texttt{FDHAContextMaker} vectorises model evaluation over all sites in the configured site collection.

For multi-site hazard-curve calculations, the framework produces one annual exceedance-rate curve per site. Output is organised so that hazard-curve quantities indexed by site---per-site annual rates of principal, distributed, and total displacement exceedance, together with the site identifier, longitude, and latitude---are written in a single results record whose row index corresponds to the input-site index. When the framework is run with a plot output, single-site plots are emitted unchanged, while multi-site plots are written as one figure per site with the site index, longitude, and latitude embedded in the figure label and the file suffix.

The calculation flow is shown in the lower half of Figure~\ref{fig:workflow}. For each fault source in the source model, the function iterates over all ruptures generated by the source's magnitude--frequency distribution and rupture-floating logic, both of which are provided by the OpenQuake Engine's source-model infrastructure. For each surface-rupturing event, the calculation proceeds as follows: (i)~an \texttt{FDHAContext} is created from the rupture; (ii)~sites are partitioned into the principal rupture zone ($r \le r_{\mathrm{threshold}}$) and the distributed rupture zone ($r > r_{\mathrm{threshold}}$), where $r_{\mathrm{threshold}}$ is an implementation default the user is expected to override based on site-specific information about the principal fault zone width and trace-mapping uncertainty; (iii)~for sites in the principal zone, the principal surface-rupture probability and principal displacement exceedance probability are evaluated using the configured models (Equation~\ref{eq:principal}); (iv)~for sites in the distributed zone, the corresponding distributed quantities are evaluated (Equation~\ref{eq:distributed}); and (v)~the contributions are accumulated as the product of the rupture's annual occurrence rate with the appropriate conditional probabilities. The total annual exceedance rate at each site is the sum of contributions from all ruptures on all fault sources, with principal and distributed contributions tracked separately so that they can be reported independently or combined in post-processing.

The \citet{Visini2025} distributed-faulting pathway is handled through a separate calculator class (\texttt{VisiniSecondaryCalculator}) because the model differs structurally from the gridding-based formulations of \citet{youngs_methodology_2003} and \citet{Petersen2011}. The Visini formulation requires classification of distributed ruptures into combinations~A, B, and~C of the SURE~2.0 rank system (see Section~\ref{sec:framework:principal_distributed}), and uses an across-strike site dimension and an along-strike site dimension rather than a single isotropic cell. The framework supports loading Rank~1.5 splay-fault traces from NRML XML files via a dedicated trace loader, enabling automatic case classification under the geologically informed decision-tree algorithm of \citet[Figure~12]{Visini2025}. As discussed in Section~\ref{sec:framework:principal_distributed}, the principal surface-rupture probability $P[\,sr \neq 0 \mid m\,]$ is not applied as a multiplicative gate in this pathway because the Visini regressions are calibrated exclusively on surface-rupturing earthquakes.


\subsection{Logic-Tree Workflow}
\label{sec:architecture:logic_tree}
The framework propagates epistemic uncertainty in PFDHA hazard estimates through two coupled logic trees: an OpenQuake-style \emph{source-model logic tree} (SMLT) that captures epistemic uncertainty in the seismic source characterization, and an \emph{FDHA model logic tree} that captures epistemic uncertainty in the choice of published PFDHA model components and selected model parameters within the four model categories. The two trees are combined as their Cartesian product: the calculation is performed once per SMLT realisation per FDHA end-branch, and combined branch weights are formed as the product of the SMLT branch weight and the FDHA-model branch weight. The two trees are independent inputs to the framework, declared in two separate NRML/XML files (Figure~\ref{fig:workflow}); this separation between configuration and the two logic trees follows OpenQuake Engine conventions \citep{Pagani2014} and allows each tree to be reused independently across sites, model selections, or alternative source models.

\paragraph{Source-model logic tree.}
\label{sec:architecture:smlt}
The INI configuration references a SMLT file under \texttt{[calculation].source\_model\_logic\_tree\_file}, and that SMLT file in turn references one or more NRML source-model XML files through the standard \texttt{sourceModel} branch set. Parsing and realisation enumeration are delegated to \texttt{openquake.hazardlib.logictree.SourceModelLogicTree}; the framework therefore honours the complete set of source-model uncertainty types recognised by the OpenQuake Engine, including alternative full \texttt{sourceModel} files, recurrence-related modifications (\texttt{bGRRelative}, \texttt{bGRAbsolute}, \texttt{abGRAbsolute}, \texttt{maxMagGRRelative}, \texttt{maxMagGRRelativeNoMoBalance}, \texttt{maxMagGRAbsolute}, \texttt{abMaxMagAbsolute}, \texttt{incrementalMFDAbsolute}, \texttt{truncatedGRFromSlipAbsolute}), geometry-related modifications (\texttt{simpleFaultGeometryAbsolute}, \texttt{simpleFaultDipRelative}, \texttt{simpleFaultDipAbsolute}, \texttt{complexFaultGeometryAbsolute}, \texttt{characteristicFaultGeometryAbsolute}, \texttt{areaSourceGeometryAbsolute}), and other supported types (\texttt{setMSRAbsolute}, \texttt{setLowerSeismDepthAbsolute}, \texttt{setUpperSeismDepthAbsolute}, \texttt{recomputeMmax}, \texttt{dummy}). When a SMLT branch carries one or more such modifications, the framework applies them through \texttt{openquake.hazardlib.lt.apply\_uncertainties}, the same engine machinery used by the OpenQuake PSHA calculator, so that the modified source set seen by the FDHA calculator is identical to what an OpenQuake PSHA run would see for the same SMLT branch.

\paragraph{FDHA model logic tree.}
The FDHA model logic tree is provided through the \texttt{[calculation].fdha\_logic\_tree\_file} entry and uses the OpenQuake NRML 0.4 logic-tree structure with FDHA-specific extensions for model-selection uncertainty. Each branch contains a model choice and an associated weight, and weights within a branch set sum to unity. Four FDHA-specific values of the \texttt{uncertaintyType} attribute identify which model category a branch set populates: \texttt{fdhaPrimarySRModel}, \texttt{fdhaPrimaryFDModel}, \texttt{fdhaSecondarySRModel}, and \texttt{fdhaSecondaryFDModel}. Branch sets also carry two structural attributes that govern how branches are combined: \texttt{applyToStyle} restricts a branch set to a specific faulting style (e.g., \texttt{reverse}), and \texttt{applyToBranches} chains a branch set onto a specified subset of branches from earlier branching levels, generating the Cartesian product of compatible branches. Per-branch model parameters that do not have config-file defaults---for example, the \texttt{cell\_size} of a Petersen-style gridding model, the displacement metric (average versus maximum) of a Youngs FD model, or the magnitude scaling relation used by a primary-SR model---are passed inside the \texttt{uncertaintyModel} block using a TOML-style inline syntax. This inline syntax allows parameter dictionaries to be associated directly with each branch without proliferating top-level configuration fields.

\paragraph{Branch chaining and model coupling.}
The \texttt{applyToBranches} attribute serves two purposes within the FDHA model logic tree. First, it enumerates the Cartesian product of branches from compatible upstream branching levels, which is the standard mechanism by which logic trees grow combinatorially. Second, it provides a controlled means of expressing model coupling constraints. For example, the \citet{Visini2025} distributed surface-rupture and distributed fault-displacement models share calibration data and parameter conventions, and were therefore developed to be used as a coupled pair; in the logic-tree file the secondary-FD branch set that selects \texttt{Visini2025SecondaryFD} can be restricted to apply only to the upstream branch that selects \texttt{Visini2025SecondarySR}, while a parallel secondary-FD branch set applies the alternative \texttt{Youngs2003SecondaryFD} to the upstream \citet{youngs_methodology_2003} and \citet{Takao2013} secondary-SR branches. This pattern ensures that physically incompatible model combinations are excluded from the Cartesian enumeration without requiring special-case logic in the calculator.

\paragraph{Branch enumeration, weight propagation, and aggregation.}
The driver enumerates the Cartesian product of SMLT realisations and FDHA end-branches. The OpenQuake NRML uncertainties recorded on that realisation are then applied before passing the source dictionary to the hazard calculator. For each combined SMLT--FDHA end-branch, a full hazard calculation is performed and a per-end-branch annual exceedance rate is stored at each site and displacement level. The combined branch weight is the product of the SMLT branch weight and the FDHA-model end-branch weight. End-branches are then aggregated to produce the mean hazard, computed as the arithmetic mean of annual rates weighted by the combined weights, and selected fractile hazard curves, here taken as the 5th, 16th, 50th, 84th, and 95th percentiles, computed across the combined-weighted end-branch ensemble. This treatment is consistent with standard PSHA logic-tree practice as implemented in the OpenQuake Engine \citep{Pagani2014, Bommer2008} and follows the principle that the logic-tree expected hazard is the arithmetic mean of branch hazard rates. The logic-tree workflow is an implementation mechanism for propagating epistemic uncertainty and does not introduce new displacement models or modify the equations, coefficients, or parameter definitions of the underlying published models.


\subsection{Configuration and Command-Line Interface}
\label{sec:architecture:config}

Calculations are specified through three input files, corresponding to the three input boxes in Figure~\ref{fig:workflow}: an INI configuration file, a NRML source-model logic-tree file, and a NRML FDHA-model logic-tree file. The INI configuration is organised into sections that follow the OpenQuake Engine \texttt{job.ini} conventions: \texttt{[general]} declares the calculation description and \texttt{calculation\_mode}; the FDHA framework introduces \texttt{fdha\_classical} as the calculation-mode value used for PFDHA, paralleling OpenQuake's \texttt{classical\_psha} mode for PSHA. The \texttt{[geometry]} section defines the site geometry and is also used to detect the calculation type: a \texttt{region} entry triggers a hazard-map calculation, while a \texttt{sites} entry (or a \texttt{sites\_csv} entry that points to a CSV file) triggers a hazard-curve calculation. The \texttt{sites} entry accepts a single \texttt{lon~lat} pair, or several whitespace-separated \texttt{lon~lat~[depth]} tokens separated by commas, allowing one or more discrete sites to be declared inline. The \texttt{sites\_csv} entry accepts the path (relative to the INI file) of a CSV file whose header includes required columns \texttt{lon} and \texttt{lat} and an optional \texttt{depth} column; longitude and latitude are validated against $(-180, 360]$ and $[-90, 90]$ in both pathways. The \texttt{[erf]} section sets the rupture-mesh and magnitude--frequency-distribution discretisation; \texttt{[calculation]} declares the paths to the two logic-tree files via \texttt{source\_model\_logic\_tree\_file} and \texttt{fdha\_logic\_tree\_file}, the target displacement levels through \texttt{displacement\_measure\_levels}, the investigation time, the return period (used only for hazard maps), and the principal-zone threshold $r_{\mathrm{threshold}}$; and \texttt{[site\_params]} sets site-specific properties such as $V_{S30}$ via \texttt{reference\_vs30\_value}. The source-model XML files referenced by the SMLT are read by the OpenQuake Engine's source-model parser without modification, so source-model XML files prepared for fault-based PSHA can be reused for PFDHA when they contain the rupture-surface information required for surface-displacement calculations. The FDHA-model logic-tree NRML file specifies the branches, weights, and chaining as described in Section~\ref{sec:architecture:logic_tree}.


A command-line interface, \texttt{fdha}, takes the INI job file as its only positional argument and dispatches through the logic-tree driver. The calculation type (hazard curve or hazard map) is auto-detected from the \texttt{[geometry]} section as described above. Per-end-branch results, per-site weighted-mean and fractile hazard curves, and (for hazard-map mode) per-end-branch HDF5 rate grids together with weighted-mean and fractile displacement maps at the configured return period are written under an output directory next to the INI file. A run-manifest JSON file records the SMLT branches and uncertainties, the FDHA end-branch selections, the combined weights, and the output file layout, supporting traceable reproduction of every reported number. Representative configuration and logic-tree files, together with the corresponding CLI invocation, are shown in Section~\ref{sec:application}.
